\section{Methodology}

\subsection{Conditional GAN architecture}

The CGAN is implemented in \code{models/gan.py} and trained in \code{train_gan.py}. It consists of:
\begin{enumerate}
\item A \textbf{Generator} that maps \((z, y)\) to a 64×64 RGB image.
\item A \textbf{Projection Discriminator} that scores whether an image is real given its class label.
\end{enumerate}

\subsubsection{Generator (conditional batch normalization)}

The generator begins with a linear projection of a noise vector \(z \in \mathbb{R}^{128}\) into a 4×4 feature map. It then applies four upsampling blocks to reach 64×64 resolution. Each block uses:
\begin{itemize}
\item Upsample ×2 (nearest-neighbor upsampling)
\item 3×3 convolution
\item \textbf{Conditional BatchNorm2d}, where \(\gamma(y)\) and \(\beta(y)\) are learned embeddings of the class label
\item ReLU activation
\end{itemize}

The final ``to RGB'' stage applies batch norm, ReLU, a 3×3 convolution to 3 channels, and \textbf{tanh} to produce outputs in \([-1, 1]\).

\subsubsection{Discriminator (projection conditioning)}

The discriminator uses a stack of convolutional residual-like blocks with average pooling for downsampling:
\begin{itemize}
\item 64×64 → 32×32 → 16×16 → 8×8 → 4×4
\end{itemize}

After the final block, it applies a leaky ReLU and \textbf{global sum pooling} over spatial dimensions to obtain a feature vector \(h(x)\). The final logit is:
\[
\ell(x, y) = \underbrace{w^\top h(x)}_{\text{unconditional}} + \underbrace{\langle e(y), h(x) \rangle}_{\text{projection}}.
\]

This logit is passed to \code{BCEWithLogitsLoss} without an explicit sigmoid layer.

\subsection{GAN training procedure}

GAN training follows a classic alternating update scheme:
\begin{itemize}
\item \textbf{One discriminator step} per batch
\item \textbf{One generator step} per batch
\end{itemize}

Losses are implemented via \code{BCEWithLogitsLoss}:
\begin{itemize}
\item Discriminator real targets: 1
\item Discriminator fake targets: 0
\item Generator targets: 1 (non-saturating objective: fool the discriminator)
\end{itemize}

\subsubsection{Loss functions (logit form)}

Let the discriminator output a raw logit \(s = D(x, y)\). Using the numerically stable \code{BCEWithLogitsLoss} corresponds to the following per-sample terms:
\begin{itemize}
\item For a real sample (target 1): \(\ell_{\text{real}} = \log(1 + \exp(-s)) = \mathrm{softplus}(-s)\)
\item For a fake sample (target 0): \(\ell_{\text{fake}} = \log(1 + \exp(s)) = \mathrm{softplus}(s)\)
\end{itemize}

With these definitions, the implemented losses correspond to:
\[
\mathcal{L}_D = \mathbb{E}_{(x,y)\sim p_{\text{data}}}\left[\mathrm{softplus}(-D(x,y))\right] + \mathbb{E}_{z,y}\left[\mathrm{softplus}(D(G(z,y),y))\right],
\]
and the non-saturating generator loss:
\[
\mathcal{L}_G = \mathbb{E}_{z,y}\left[\mathrm{softplus}(-D(G(z,y),y))\right].
\]

This formulation matches the ``target=1'' generator objective implemented in the training code.

\subsubsection{Update details (as implemented)}

The discriminator and generator updates are not purely symmetric; the specific conditioning choices matter:
\begin{itemize}
\item \textbf{Discriminator step.} Fake images are generated using the \emph{real batch labels} \(y_{\text{real}}\). This forces the discriminator to evaluate real-vs-fake \emph{at the same conditional label} for each sample.
\item \textbf{Generator step.} Labels are sampled uniformly at random. This encourages the generator to model all classes and not only mimic the label distribution in a particular batch.
\end{itemize}

In pseudocode (per minibatch):
\begin{enumerate}
\item Sample a real minibatch \((x_{\text{real}}, y_{\text{real}})\).
\item Generate \(x_{\text{fake}} = G(z, y_{\text{real}})\) for discriminator training.
\item Update \(D\) to classify \(x_{\text{real}}\) as real and \(x_{\text{fake}}\) as fake.
\item Sample random labels \(y \sim \mathrm{Unif}(\{1,\dots,K\})\), generate \(x_{\text{fake}} = G(z, y)\).
\item Update \(G\) to make \(D(x_{\text{fake}}, y)\) predict ``real''.
\end{enumerate}

Optimization uses Adam with betas \((0.5, 0.999)\), a common choice in DCGAN-style training. To reduce instability, \textbf{gradient clipping} is applied to both networks with \code{max_norm=1.0}.

The training loop periodically:
\begin{itemize}
\item Saves sample grids (\code{samples_epochXXXX.png}) for qualitative monitoring
\item Saves checkpoints (\code{ckpt_epochXXXX.pt})
\item Computes FID every \code{fid_every} epochs on a fixed evaluation split
\end{itemize}

Qualitative sample grids are generated from a fixed set of latent vectors and labels (8 samples per class). Using fixed inputs reduces visual ``jitter'' and makes it easier to attribute changes in sample quality to training progress rather than random sampling variation.

\subsubsection{Checkpoint selection}

Whenever an evaluated checkpoint improves the best recorded FID, it is saved as \code{runs/gan/checkpoints/best_fid.pt}. This report uses that selection rule for synthetic pool generation, ensuring that the synthetic dataset corresponds to the best intrinsic metric observed during training (subject to the limitations of FID discussed earlier).

\subsection{FID computation details (as implemented)}

FID is computed in \code{train_gan.py} using:
\begin{itemize}
\item Feature extractor: InceptionV3 with ImageNet weights (\code{IMAGENET1K_V1})
\item Output: 2048-d ``pool3'' features (implemented by replacing the final FC with identity)
\item Preprocessing:
\begin{itemize}
\item Generator outputs in \([-1,1]\) are rescaled to \([0,1]\)
\item Then normalized with ImageNet mean/std
\item Resized to 299×299 before feeding into InceptionV3
\end{itemize}
\end{itemize}

The covariance matrix square root \((\Sigma_r \Sigma_g)^{1/2}\) is computed using \code{scipy.linalg.sqrtm}. Numerical errors can introduce small imaginary components; the implementation discards the imaginary part when it occurs, which is a common practical handling strategy in FID implementations.

The mean and covariance are estimated from finite samples using NumPy (\code{mean} and \code{np.cov}). By default, \code{np.cov} uses the sample covariance normalization (approximately \(1/(N-1)\)), which is standard in many FID implementations. However, when \(N\) is small (as in the 477-image validation split), covariance estimates can be noisy, contributing to non-monotonic FID curves.

\textbf{Important practical note.} The configuration requests \code{fid_n_samples=2048}, but the validation split contains only \textbf{477} images. Therefore, each FID evaluation uses \textbf{477 real + 477 generated} images in this setup. This can increase variance and means the absolute FID values should be interpreted cautiously; the primary signal is the \emph{trend} and relative comparisons across epochs.

\subsection{Synthetic pool generation}

Synthetic images are generated using \code{scripts/generate_synth.py} by loading a trained generator checkpoint and sampling \(z \sim \mathcal{N}(0, I)\) with class labels \(y \in \{0,1,2\}\). The output pool is stored in:
\begin{itemize}
\item \code{data_synth/{apple,banana,orange}/...png}
\end{itemize}

For the downstream experiments, the synthetic pool is balanced across classes.

In the current repository state, the synthetic pool contains \textbf{1300 images per class} (3900 total), matching the full size of the real training split. Synthetic images are saved as PNG files with normalization parameters chosen so that pixel values correspond to the generator's \([-1,1]\) output range when visualized.

\subsection{Downstream classifier}

The classifier (\code{models/classifier.py}) is a small convolutional neural network (FruitCNN) with:
\begin{itemize}
\item Repeated Conv → BatchNorm → ReLU layers
\item MaxPool downsampling
\item Dropout for regularization
\item A 256-unit fully connected layer and a final linear classifier
\end{itemize}

Training uses cross-entropy loss, Adam optimizer, and a cosine annealing learning rate schedule.

The classifier is intentionally modest in capacity to make data-size effects visible and to reduce the chance that results are dominated by overparameterization. Dropout is applied both in the convolutional stack (Dropout2d) and in the fully connected head to regularize training across all scenarios.

\subsection{Reproducibility considerations}

The repository is designed to be reproducible at a practical level (fixed seeds and deterministic subsampling), but strict determinism is not guaranteed:
\begin{itemize}
\item \textbf{Random seeds.} \code{train_gan.py} seeds Python \code{random}, NumPy, and PyTorch. \code{train_classifier.py} seeds Python \code{random} and PyTorch. The deterministic subsampling by class uses \code{seed=42}.
\item \textbf{Data loading.} Multi-worker data loading can introduce subtle nondeterminism unless worker seeds and deterministic algorithms are enforced.
\item \textbf{Hardware backend.} The default device is Apple Silicon MPS when available. Some backend operations can be nondeterministic across runs or across software versions.
\end{itemize}

For a thesis-grade evaluation, the recommended practice is to run multiple seeds and report aggregate statistics (mean ± standard deviation) rather than relying on exact bitwise reproduction.
